% Part: first-order-logic
% Chapter: introduction
% Section: semantic-notions

\documentclass[../../../include/open-logic-section]{subfiles}

\begin{document}

\olfileid{fol}{int}{sem}

\olsection{Gagasan Semantis}

Ing ndhuwur wis kasebut yen nalika nimbang apa $\Sat{M}{!A}[s]$
lumaku, kita (supaya gampang) marengake $s$ menehi nilai marang kabeh
!!{variable}s, nanging mung nilai kang diwenehake marang !!{variable}s
ing~$!A$ kang digunakake. Malah, mung nilai kanggo variabel
\emph{bebas} ing~$!A$ kang nduweni pengaruh. Mesthi wae, amarga kita
ngati-ati, kasunyatan iki bakal kita buktekake. Amarga !!{sentence}s
ora nduweni variabel bebas, $s$ ora nduweni pengaruh babar pisan
marang dicukupi utawa orane ukara kasebut ing !!a{structure}. Mula,
nalika $!A$ iku !!a{sentence}, kita bisa netepake $\Sat{M}{!A}$ kanthi
teges ``$\Sat{M}{!A}[s]$ kanggo saben~$s$,'' kang nyatane bener yen lan
mung yen $\Sat{M}{!A}[s]$ kanggo paling ora siji~$s$. Kita kudu
ngenalake pangaji !!{variable} supaya bisa menehi definisi relasi
nyukupi kang lumaku kanggo !!{formula}s, nanging kanggo !!{sentence}s,
relasi nyukupi ora gumantung marang pangaji !!{variable}.

Sawise nduweni definisi ``$\Sat{M}{!A}$,'' kita ngerti tegesing
``kasus'' lan ``bener ing'' tumrap !!{sentence}s ing logika sepisanan.
Adhedhasar definisi $\Sat{M}{!A}$ kanggo !!{sentence}s, kita banjur bisa
nemtokake gagasan semantis dhasar yaiku validitas, akibat semantis, lan
kabisane dicukupi. Sawijining ukara iku valid, $\Entails !A$, yen saben
!!{structure} nyukupi ukara kasebut. Ukara iku dadi akibat saka
sawijining himpunan !!{sentence}s, $\Gamma \Entails !A$, yen saben
!!{structure} kang nyukupi kabeh !!{sentence}s ing~$\Gamma$ uga
nyukupi~$!A$. Lan sawijining himpunan !!{sentence}s bisa dicukupi yen
ana !!{structure} kang nyukupi kabeh !!{sentence}s ing himpunan
kasebut bebarengan.

Amarga !!{formula}s ditemtokake kanthi induktif, lan relasi nyukupi
sabanjure ditemtokake kanthi induksi ing struktur !!{formula}s, kita
bisa nggunakake induksi kanggo mbuktekake sipat-sipat semantik kita
lan nggandhengake gagasan-gagasan semantis kang wis ditemtokake. Kita
bakal nglumpukake lan mbuktekake sawetara sipat kasebut, saperangan
amarga saben sipat kasebut narik kawigaten dhewe, nanging utamane
amarga akeh kang bakal migunani nalika kita nerusake panliten ngenani
relasi antarane semantik lan sistem !!{derivation}. Kanggo nindakake
iki, kita uga kudu nemtokake (kanthi trep, yaiku kanthi induksi)
sawetara gagasan lan operasi sintaksis kang durung kasebut.

\end{document}
